% Act 2 example 1 of 5. Shared simple template for short/12--16: three boxes (Shared state -> Alternatives -> Check),
% one big result line, one claim with gold rule chips. The \providecommand block below is copied verbatim into
% all five files; the first one loaded defines it, so every example draws the same geometry. Edit all five together.
% #1 shared state, #2 alternatives, #3 check, #4 result line (use \exnum{} for the big numbers).
\providecommand{\exnum}[1]{{\fontsize{20}{24}\selectfont\bfseries #1}}
\providecommand{\rulechip}[1]{\tikz[baseline=(x.base)]\node[fill=Gold,text=white,rounded corners=2pt,inner xsep=3.5pt,inner ysep=2.4pt,font=\fontsize{11}{13}\selectfont\bfseries](x){#1};}
\providecommand{\exampleflow}[4]{%
\noindent\begin{tikzpicture}[x=1cm,y=1cm,
  hd/.style={anchor=south,font=\fontsize{11.5}{13}\selectfont\bfseries,text=Teal},
  bx/.style={text width=3.9cm,inner xsep=.15cm,minimum height=1.45cm,font=\fontsize{13.5}{16}\selectfont\bfseries},
  stack/.style={draw=Gold,fill=Sand,line width=.9pt},
  res/.style={anchor=west,inner sep=0pt,font=\fontsize{15}{18}\selectfont},
  every node/.style={align=center}]
\useasboundingbox (0,.62) rectangle (14.4,-2.85);
\node[hd] at (2.12,.15) {Shared state};
\node[hd] at (7.10,.15) {Alternatives};
\node[hd] at (12.28,.15) {Check};
\node[state,bx] (s) at (2.12,-.8) {#1};
\node[candidate,bx] (a) at (7.10,-.8) {#2};
\begin{scope}[on background layer]
\draw[stack] ($(a.south west)+(.2,.2)$) rectangle ($(a.north east)+(.2,.2)$);
\draw[stack] ($(a.south west)+(.1,.1)$) rectangle ($(a.north east)+(.1,.1)$);
\end{scope}
\node[evidence,bx] (c) at (12.28,-.8) {#3};
\draw[flow] (s.east)--(a.west);
\draw[flow] ($(a.east)+(.2,0)$)--(c.west);
\node[res] at (0,-2.35) {#4};
\end{tikzpicture}\par}
\begin{frame}[t]{\ev{\tagO}\surf{SOFTWARE SECURITY}Reuse makes fuzzing fast; only checks turn crashes into findings.\strut}
\vspace{.2cm}
\exampleflow{Target\\initialized once}{Changed\\inputs}{Reproduce,\\deduplicate}
{AFL++\enspace\exnum{10--20$\times$} faster\qquad CyberGym\enspace\exnum{56\,$\to$\,22} confirmed}
\claim{Set up once; count each bug once.\hfill\rulechip{COMPUTE ONCE}\enspace\rulechip{COUNT ONCE}}
\sources{Two separate sources. \href{https://github.com/AFLplusplus/AFLplusplus/blob/stable/instrumentation/README.persistent_mode.md}{AFL++ persistent-mode docs}: typical gain vs a fork per input. \href{https://www.cybergym.io/cybergym/}{CyberGym report}: OpenHands + GPT-5 wrote and ran inputs on 431 projects (latest-code campaign).}
\qadetail{Two separate sources: AFL++ documentation supplies the speed figure; CyberGym is a separate agent study that does not claim AFL++ or persistent-mode reuse. AFL++: the documented typical 10--20x gain is for persistent mode, which avoids a process fork for every input; it must not be attributed solely to the deferred forkserver. The forkserver defers cloning until initialization is complete; a persistent child then runs the harness repeatedly, resetting mutable state between inputs. About 1,000 iterations before restarting is a recommended starting point, not a fixed setting. Gains depend on the target and reset correctness. A newly executed input is not automatically an independent observation, and a crash is not automatically an exploitable vulnerability. CyberGym latest-code campaign (report checked 4 October 2026): OpenHands with GPT-5 on 431 latest OSS-Fuzz projects and 1,748 runnable executables produced 56 crashes; the authors reproduced, deduplicated and investigated them, confirming 22 distinct previously unknown vulnerabilities. Projects, targets, crashes and validated bugs are different units; the funnel is not a disjoint 22/56 success probability, and exploitability or patch status is not implied. This is separate from CyberGym's historical 1,507-vulnerability benchmark across 188 projects. The separate GPT-4.1 campaign's counts overlap and cannot be added. Neither result measures sandbox forking or our own security campaign.}
% Transition: from the rules map into the first example; ends announcing AI research.
\note{[1:00] First example: software security. A fuzzer runs a program on many changed inputs, hunting for crashes. AFL-plus-plus initializes the target once, then reuses one child process across many inputs. Its documentation reports typical speedups of ten to twenty times. It skips a fork for every input. That works only if state resets correctly. That is rule one, compute once. But speed only produces leads, and leads need checks. Berkeley's CyberGym is a separate study. There, an AI agent wrote test inputs for 431 projects. It produced fifty-six crashes. The researchers reproduced them and removed duplicates, leaving twenty-two distinct new vulnerabilities. That is rule two, count once. Second example: AI research.}
\end{frame}
